\documentclass[10pt,a4paper]{amsart}
\usepackage[margin=19mm]{geometry}
\usepackage{amsmath,amssymb,mathtools}
\usepackage[scr=rsfs,cal=euler]{mathalfa}
\usepackage{xfrac}
\usepackage[unicode,hidelinks,bookmarks=false]{hyperref}
\newcommand{\ie}{i.e.\ }
\newcommand{\cf}{cf.\ }
\newcommand{\resp}{respectively\ }
\newcommand{\Subsection}{\subsection}
\providecommand{\nobreakdash}{\nobreak\hbox{-}}
\setlength{\emergencystretch}{2em}
\multlinegap0pt
\usepackage{bm}
\usepackage{bbm}
\usepackage{dsfont}
\usepackage{xspace}
\usepackage{cancel}
\usepackage{mathtools}
\usepackage{amsthm}
\makeatletter
\@ifundefined{thm}{\newtheorem{thm}{Thm}}{}
\@ifundefined{lem}{\newtheorem{lem}[thm]{Lemma}}{}
\makeatother
\makeatletter
\expandafter\ifx\csname example\endcsname\relax
\newenvironment{example}[1][Example]{\begin{trivlist}\item[\hskip \labelsep {\bfseries #1}]}{\end{trivlist}}
\fi
\makeatother
\title[Minimal Unimodal Decomposition is NP-Hard on Graphs]{Minimal Unimodal Decomposition is NP-Hard on Graphs}
\author{Mishal Assif P K, Yuliy Baryshnikov}
\date{Source: arXiv:2510.05944v1 (CC BY 4.0)}
\begin{document}
\maketitle

\noindent\emph{Source and licence.} Minimal Unimodal Decomposition is NP-Hard on Graphs. Version arXiv:2510.05944v1 (CC BY 4.0), distributed under the Creative Commons Attribution 4.0 International licence, as recorded in the arXiv licence metadata for this exact version and shown on the abstract page https://arxiv.org/abs/2510.05944v1. Full rights notice: licenses/arxiv-2510.05944-CC-BY-4.0.txt.\par\smallskip

\noindent\textit{Editorial scope.} Counted unit: the Lem whose source label is \texttt{lem:umd} in the article (source0.tex lines 679--682, source label \texttt{lem:umd}), delivered together with the article's own complete original proof of it (source0.tex lines 683--740, one \texttt{proof} environment of 58 lines). No mathematics is written, repaired or re-derived: statement and proof are the article's own text line for line. Required context retained uncounted: none. Every definition, display and label of those lines is retained unchanged. Omitted from this package and not redistributed: 1--678, 741--939. Nothing inside a retained excerpt is omitted or altered: every retained body is delivered exactly as the article's lines are written, so no omission receipt of any kind accompanies this package. Citations: the retained excerpts carry no citation command at all, so no bibliography entry of the article is transcribed and no reference is redistributed. Reference adaptation: none was needed and none was made; every reference inside the delivered unit resolves inside it, so no unresolved reference or citation is produced. Printed slips kept literal: no printed word, symbol, hypothesis or number of the counted statement or of its proof was altered. Layout and class: the article's own class is \texttt{article} with packages including geometry, eulervm, tgpagella, fontenc, amsmath, amssymb, amsthm, cite, color, graphicx, bm, mathtools, algorithm, algpseudocode; the wrapper is the lane's pinned \texttt{amsart} editorial wrapper at 10pt on A4, which restates the theorem-like environments the retained text needs and adds the article's own macro definitions that the retained text needs (none), with extra wrapper packages (bm, bbm, dsfont, xspace, cancel, mathtools). The delivered numbering is the unit's own. Assets: no retained span contains a figure environment, a graphics-inclusion command, a PDF-page inclusion, a table or a TikZ picture; no image file is copied into this package, the package declares no asset dependency and no image file is redistributed with it. Licence: version arXiv:2510.05944v1 of this article is distributed under the Creative Commons Attribution 4.0 International licence, as recorded in the arXiv licence metadata for this exact version and shown on the abstract page https://arxiv.org/abs/2510.05944v1. A full rights notice is delivered at licenses/arxiv-2510.05944-CC-BY-4.0.txt. Characters: every retained excerpt is pure ASCII, so the delivered engine, XeLaTeX with its default Latin Modern text font, reports no missing character.
\begin{lem}
\label{lem:umd}
    If $\tilde{f}$ has a unimodal decomposition with $k$ components, then $G$ has a vertex cover with at most $k$ components.
\end{lem}

\begin{proof}
Now assume that $\tilde{f}$ has a unimodal decomposition with $k$ components $F = \{f_1, \ldots, f_k\}$. We pick at most one mode from each component to form a set of vertices in $G$. If the only mode of a component function is a vertex in $\tilde{G} - G$, we ignore this mode. If a component has multiple modes, with at least one mode in $G$, we pick any one of those modes in $G$. We will prove that the set of vertices in $G$ we obtain in this manner forms a vertex cover with at most $k$ elements. We will prove by contradiction. Let us assume that there are two vertices in $G$, $v_1$ and $v_2$ such that both the vertices are modes of none of the components. We will show that these two vertices can't be adjacent in $G$. Let $v$ be a vertex that is not the mode of any function in $F$, and $E_v = \{e_1, \ldots, e_{\deg(v)}\}$ be the edges hitting $v$ in $G$. Define $F^{v}_{e}$ to be the set of functions
\[
	F^{v}_{e} = \{ f_i \in F \ | f_i(e)  \geq f_i(v) \geq f_i(e^{'}) \ \forall e^{'} \neq e \}.
\]
The union of the collection $\{ F^{v}_{e} \}_{e \in E_v}$ should equal $F$, since $v$ is not the mode of any function in $F$. In addition, 
\[
	\sum_{f \in F^{v}_{e}} f(v) \leq \sum_{f \in F^{v}_{e}} f(e) \leq 1.
\]
which means 
\begin{equation}
\label{impsum1}
	\deg(v) = \sum_{f \in F} f(v) = \sum_{e \in E_v} \left( \sum_{f \in F^{v}_{e}} f(v) \right) \leq \sum_{e \in E_v} 1 = \deg(v).
\end{equation}
This implies that for each $e \in E_v$,
\[
	\sum_{f \in F^{v}_{e}} f(v) = 1.
\]
In addition, if $e_1 \neq e_2 \in E_v$ and $f \in F^{v}_{e_1}$ and $f(e_1) > 0$, then $f(e_2) = 0$. If $f(e_2) > 0$, there are two possibilities:
\begin{itemize}
	\item Either $f(e_2) <  f(e_1)$. In this case, $f \notin F^{v}_{e_2}$ but $f(e_2) = a > 0$, which means
	\[
		\sum_{g \in F^{v}_{e_2}} g(e_2) = 1 - \sum_{g \in F-F^{v}_{e_2}} g(e_2) \leq 1-f(e_2) < 1, \quad \text{since $f \notin F^{v}_{e_2}$,}
	\] 
	which further implies
	\[
		\deg(v) = \sum_{g \in F} g(v) \leq \sum_{e \in E_v} \left( \sum_{g \in F^{v}_{e}} g(v) \right)  < \deg(v)
	\]
	giving a contradiction.
	\item Or $f(e_2) = f(v) = f(e_1) > 0$. In this case $f \in F^v_{e_2} \cap F^v_{e_1}$, which means
	\[
		\deg(v) = \sum_{g \in F} g(v) \leq \sum_{e \in E_v} \left( \sum_{g \in F^{v}_{e}} g(v) \right)  - f(v) \leq \deg(v) - f(v) < \deg(v),
	\]
	where the first inequality follows from the fact that when $F$ is written as $\cup_{e \in E_v} F^{v}_{e}$, the function $f$ appears in both $F^v_{e_2}$ and $F^v_{e_1}$, which means it is counted twice and can be subtracted once from the sum on the right side of the inequality. This again gives a contradiction.
\end{itemize}
To summarize, the collection of functions $\{ F^{v}_{e} \}_{e \in E_v}$ satisfy the following two properties:
\begin{equation}
\label{prop1-mode-free}
\cup_{e \in E_v} F^{v}_{e} = F,
\end{equation}
\begin{equation}
\label{prop3-mode-free}
\sum_{f \in F^{v}_{e}} f(v) = 1,
\end{equation}
and
\begin{equation}
\label{prop2-mode-free}
f \in F^{v}_{e} \implies f(e^{'}) = 0 \ \text{ for all $e^{'} \in E_v - \{e\}$}.
\end{equation}

Now suppose $v_1$ and $v_2$ are vertices in $G$ with edge $e_{1,2}$ hitting both vertices. If both $v_1$ and $v_2$ are not the modes of any functions in $F$, then the set of functions $F^{v_1}_{e_{1,2}}$ = $F^{v_2}_{e_{1,2}} $. This is true as $f(e_{1,2}) \geq f(v_1)$ implies $f(e_{1,2}) \geq f(v_2)$. To see this, observe that if $f(v_2) > f(e_{1,2})$ then $f \in F^{v_2}_{e^{'}}$ for a different edge $e^{'}$ and by property \eqref{prop2-mode-free} $f(e_{1,2})$ would be zero. 
In addition, any function in the set $F^{v_1}_{e_{1,2}}$ must take value zero on any vertex in $\tilde{G}$ other than $v_1, e_{1,2}, v_2$, also due to property \eqref{prop2-mode-free} and unimodality of $f$. Now we can see that
\begin{equation}
\label{finalprop}
	\sum_{f \in F^{v_1}_{e_{1,2}}} f(v_1) = \sum_{f \in F^{v_1}_{e_{1,2}}} f(v_2) = \sum_{f \in F^{v_1}_{e_{1,2}}} f(e_{1,2}) = 1, \text{ \ and \ } \sum_{f \in F^{v_1}_{e_{1,2}}} f(v) = 0 \ \text{ for all other $v \in \tilde{G}$} 
\end{equation}
This means we can remove the set of functions $F^{v_1}_{e_{1,2}}$ from $F$ and replace it with a single function $\sum_{f \in F^{v_1}_{e_{1,2}}} f$, a  unimodal function as it takes the value $1$ on a path containing three vertices and value $0$ on all other vertices. This modification can be done in polynomial time, by scanning through each vertex in $\tilde{G}-G$ to check if it satisfies \eqref{finalprop}, which can take at most $O((E + V)^3)$ time. After this modification, both $v_1$ and $v_2$ are modes of the new function $\sum_{f \in F^{v_1}_{e_{1,2}}} f$, which is a contradiction.
\end{proof}

\end{document}
